双路径单光子态为 \(|\psi_\phi\rangle=(|0\rangle+e^{i\phi}|1\rangle)/\sqrt2\)，未知参数是相位 \(\phi\)。每个光子独立存活概率 \(0<\eta\le1\)；损失标记与双路径子空间正交且不包含相位。实验重复 \(\nu\) 次，采用局部无偏估计及真实入射光子数作资源账本。
(1) 由纯态 \(F_Q=4(\langle\psi'|\psi'\rangle-|\langle\psi|\psi'\rangle|^2)\) 求量子Fisher信息和Cramér-Rao界。
(2) 在工作点构造达到该QFI的两结果投影测量，计算概率和经典Fisher信息，讨论相位偏置。
(3) 对单光子损失旗标混合态求可用QFI，解释只保留成功试验而不计失败为何误报精度。
(4) 对NOON态 \((|N,0\rangle+e^{iN\phi}|0,N\rangle)/\sqrt2\) 求有无独立损耗的QFI，并在以入射光子为代价的归一化下优化 \(N\) 的尺度。
